\section{From public API to executable control path}\label{sec:C22:codepaths}
A concrete audit can start with three interfaces rather than an entire framework taxonomy. In the inspected LangGraph source tree, the state-graph implementation resides in the graph state module and exposes compilation into an executable graph. In the OpenAI Agents Python SDK, the run module exposes the Runner control surface. The Google ADK source exposes runner infrastructure for application execution. These are identifiable inspection points, not assertions that the corresponding modules implement every guarantee discussed above. Their current source files were reachable at the cutoff; experiment revisions still require immutable pinning~\cite{executionproduct001,executionproduct002,executionproduct003}.

For the graph case, trace compilation into checkpoint use. Identify where thread identity enters configuration, which state updates are reduced together, and which node outputs are persisted before successors execute. Then place an effectful tool call inside one node and inject a crash immediately after the service accepts it. If resumption reexecutes the node, the application must supply idempotency or reconciliation. This test distinguishes graph-state durability from external-effect uniqueness without requiring a language model in the fixture.

For the Runner case, trace model output through tool-call decoding, schema validation, execution, and subsequent model input. Capture the exact error object returned for an invalid argument. Then inspect whether a handoff changes the active instructions, tools, model settings, or state ownership. The experiment should preserve the pre-handoff authority envelope and verify that delegation cannot expand it. A handoff that routes a conversation successfully can still be incorrect if it transfers more context or credentials than the receiving specialist requires.

Current source inspection supplies narrower testable claims. LangGraph's pinned \path{DeltaChannel} constructor sets \texttt{snapshot\_frequency=100}; this supports the checkpoint-layout claim without asserting semantic preservation by compaction~\cite{cumb2026langgraphcode}. The Agents SDK's pinned \path{RunState} uses schema 1.21 and rejects ambiguous legacy agent identities; its documentation warns that neither context overrides nor strict context deserialization establishes snapshot integrity~\cite{cumb2026agentsstate}. The required fixture therefore changes a duplicate-agent graph and verifies a safe rejection, rather than relabeling an old checkpoint until loading succeeds. Serialization compatibility and authenticity are separate checks.

For an ADK application, inspect how an event updates session state and how workflow execution consumes that state. Duplicate one event identifier and delay another. The expected behavior must come from the selected application's contract, because a generic toolkit may legitimately expose several scheduling modes. Evaluate whether the serialized event contains sufficient identity to diagnose the difference. A toolkit's flexibility increases the number of legal configurations; it does not eliminate the need to state which one produced a result.

For example, the pinned ADK \path{InMemorySessionService.append_event} skips partial events and suppresses an equal re-delivery only after matching its event identifier. Distinct event objects with the same identifier but different fields are retained. Its scope is an in-memory session implementation, not restart durability or universal event-ID uniqueness across backends~\cite{cumb2026adksession}. A discriminating fixture must therefore test both identical re-delivery and conflicting reuse, and must inspect the selected persistent backend independently. An identifier's presence in a serialized object is insufficient evidence that every consumer enforces its intended uniqueness.

These inspections produce a claim-to-code ledger. Each row records the public claim, source revision, module and callable, relevant configuration, injected condition, expected invariant, and observed outcome. Until the fixture has run, the outcome field remains unmeasured. This restraint prevents source inspection from becoming an invented reproduction. When the experiment does run, a discrepancy should lead to a narrower conclusion: perhaps a configured checkpointer preserves local state, while the selected tool adapter repeats effects. Blaming or endorsing the entire framework would discard the evidence about where the behavior actually arose.

The ledger also identifies work needed for upgrades. If a release changes the callable responsible for event reduction or tool retries, rerun the fixtures attached to those claims. An unrelated documentation edit need not trigger all experiments. This dependency-based approach makes validation scalable while retaining scientific traceability. It is especially useful for agent systems because a small change to error serialization can alter the model's next action even when the application's visible API remains compatible.
